\documentclass[aspectratio=169]{beamer}

% Tema UFS --- Universidade Federal de Sergipe
\usetheme{UFS}

% Pacotes base
\usepackage[utf8]{inputenc}
\usepackage[portuguese]{babel}
\usepackage{amsmath,amssymb}
\usepackage{graphicx}
\usepackage{booktabs}
\usepackage{xcolor}

% TikZ e bibliotecas
\usepackage{tikz}
\usetikzlibrary{shapes.geometric, arrows.meta, positioning, matrix, fit,
  backgrounds, decorations.pathreplacing, shadows, patterns, calc}

% Listagens --- paleta VS Code Light+
\usepackage{listings}
\definecolor{bgcolor}{HTML}{FFFFFF}
\definecolor{linecolor}{HTML}{DDDDDD}
\definecolor{numcolor}{HTML}{999999}
\definecolor{kwcolor}{HTML}{0000FF}
\definecolor{strcolor}{HTML}{A31515}
\definecolor{cmtcolor}{HTML}{008000}

% listings nao conhece JavaScript/JSX: definimos a linguagem
\lstdefinelanguage{JavaScript}{
  morekeywords={const,let,var,function,return,import,from,export,default,
    if,else,for,while,new,class,extends,super,this,typeof,null,undefined,
    true,false,async,await,of,in,delete},
  morecomment=[l]{//},
  morecomment=[s]{/*}{*/},
  morestring=[b]',
  morestring=[b]",
  sensitive=true,
}

\lstdefinestyle{ufsStyle}{
  backgroundcolor=\color{bgcolor},
  basicstyle=\ttfamily\footnotesize\color{black},
  keywordstyle=\color{kwcolor}\bfseries,
  stringstyle=\color{strcolor},
  commentstyle=\color{cmtcolor}\itshape,
  numberstyle=\tiny\color{numcolor},
  numbers=left, numbersep=12pt,
  xleftmargin=20pt, framexleftmargin=20pt,
  breaklines=true, tabsize=2,
  keepspaces=true, showspaces=false,
  showstringspaces=false, showtabs=false,
  frame=single, rulecolor=\color{linecolor},
  framesep=4pt, captionpos=b, upquote=true,
  columns=fullflexible,
  literate=%
    {á}{{\'a}}1 {é}{{\'e}}1 {í}{{\'i}}1 {ó}{{\'o}}1 {ú}{{\'u}}1
    {â}{{\^a}}1 {ê}{{\^e}}1 {ô}{{\^o}}1 {ã}{{\~a}}1 {õ}{{\~o}}1
    {ç}{{\c c}}1 {Á}{{\'A}}1 {É}{{\'E}}1 {Í}{{\'I}}1 {Ó}{{\'O}}1
    {Ú}{{\'U}}1 {Ç}{{\c C}}1 {à}{{\`a}}1 {`}{{\`{}}}1,
}
\lstset{style=ufsStyle, language=JavaScript}

\lstdefinestyle{js}{
  style=ufsStyle, language=JavaScript, numbers=none,
  basicstyle=\ttfamily\scriptsize\color{black},
  aboveskip=4pt, belowskip=1pt,
}
\lstdefinestyle{jsx}{
  style=ufsStyle, language=JavaScript, numbers=none,
  basicstyle=\ttfamily\tiny\color{black},
  aboveskip=3pt, belowskip=1pt,
}

% --- Estilos TikZ reutilizados ---
\tikzset{
  cx/.style={rectangle, draw=UFSBlue, fill=UFSBlue!8, rounded corners,
    minimum width=2.8cm, minimum height=0.8cm, font=\scriptsize, align=center},
  cxr/.style={rectangle, draw=UFSRed, fill=UFSRed!8, rounded corners,
    minimum width=2.8cm, minimum height=0.8cm, font=\scriptsize, align=center},
  cxy/.style={rectangle, draw=UFSYellow!80!black, fill=UFSYellow!20,
    rounded corners, minimum width=2.8cm, minimum height=0.8cm,
    font=\scriptsize, align=center},
  cxg/.style={rectangle, draw=UFSGray, fill=black!4, rounded corners,
    minimum width=2.8cm, minimum height=0.8cm, font=\scriptsize, align=center},
  cel/.style={rectangle, draw=UFSGray, fill=white, rounded corners=3pt,
    minimum width=2.6cm, minimum height=0.55cm, font=\tiny, align=center},
  seta/.style={->, thick, UFSBlue, >={Stealth[length=4pt]}},
  setar/.style={->, thick, UFSRed, >={Stealth[length=4pt]}},
  setag/.style={->, thick, UFSGray, >={Stealth[length=4pt]}},
  rot/.style={font=\tiny, text=UFSGray},
  item/.style={rectangle, draw=UFSBlue, fill=UFSBlue!10, rounded corners,
    minimum width=10.2cm, minimum height=0.68cm, font=\small},
}

% Referencia da fonte no rodape do frame
\newcommand{\fonte}[1]{\vfill\hfill{\scriptsize\textcolor{UFSGray}{Fonte: #1}}}

% --- Metadados ---
\title{Fundamentos de React Native}
\subtitle{E as suas primeiras telas rodando no celular}
\author{Prof.\ André Yoshiaki Kashiwabara}
\institute{DCOMP -- Universidade Federal de Sergipe\\
\texttt{andreyoshiaki@dcomp.ufs.br}}
\date{}

\begin{document}

\begin{frame}[plain]
  \titlepage
\end{frame}

% ==================================================================
\begin{frame}{O que você vai aprender hoje}
  \begin{center}
  \begin{tikzpicture}
    \node[item] (t1) at (0, 0)    {Como um aplicativo chega até a tela do celular};
    \node[item] (t2) at (0,-0.90) {O mínimo de JavaScript e React para começar};
    \node[item] (t3) at (0,-1.80) {Sua primeira tela rodando no \emph{seu} aparelho};
    \node[item] (t4) at (0,-2.70) {Componentes, estilos e a organização do layout};
    \node[item] (t5) at (0,-3.60) {Telas que reagem ao toque e ao que você digita};
    \draw[seta] (t1) -- (t2);
    \draw[seta] (t2) -- (t3);
    \draw[seta] (t3) -- (t4);
    \draw[seta] (t4) -- (t5);
  \end{tikzpicture}
  \end{center}
\end{frame}

% ==================================================================
\begin{frame}{Onde a aula termina}
  \begin{columns}[T]
    \begin{column}{0.54\textwidth}
      \textbf{A meta de hoje}\\[0.3em]
      Daqui a 150 minutos o seu celular vai estar rodando um
      \emph{cartão de perfil} escrito por você: foto, nome, bio editável,
      etiquetas de habilidades, um botão de curtir que conta e um botão
      de seguir que muda de cor.

      \medskip
      \textbf{O que você precisa}\\[0.3em]
      Um navegador, o celular com o aplicativo \textbf{Expo Go} instalado
      e a rede da universidade. Nada para compilar, nada para instalar
      no computador.
    \end{column}
    \begin{column}{0.44\textwidth}
      \begin{center}
      \begin{tikzpicture}[scale=0.95, transform shape]
        \draw[rounded corners=6pt, draw=UFSGray, fill=black!3]
          (0,0) rectangle (4.1,4.3);
        \draw[fill=UFSBlue!25, draw=UFSBlue] (0.75,3.55) circle (0.42);
        \node[font=\scriptsize\bfseries, anchor=west] at (1.35,3.75) {Ada Lovelace};
        \node[font=\tiny, text=UFSGray, anchor=west] at (1.35,3.42) {@ada · 128 seguidores};
        \node[font=\tiny, text=UFSGray, align=left, anchor=north west]
          at (0.35,2.95) {Estudante de Programação B.\\Gosto de algoritmos e de café.};
        \node[cel, minimum width=1.0cm, minimum height=0.35cm] at (0.85,2.05) {JavaScript};
        \node[cel, minimum width=1.1cm, minimum height=0.35cm] at (2.1,2.05) {React Native};
        \node[cel, minimum width=0.8cm, minimum height=0.35cm] at (3.3,2.05) {Flexbox};
        \draw[rounded corners=4pt, fill=UFSGray!15, draw=UFSGray]
          (0.35,1.25) rectangle (3.75,1.8);
        \node[font=\tiny, text=UFSGray] at (2.05,1.52) {campo de edição da bio};
        \draw[rounded corners=4pt, fill=UFSRed, draw=UFSRed] (0.35,0.5) rectangle (1.95,1.0);
        \node[font=\tiny, text=white] at (1.15,0.75) {7 curtidas};
        \draw[rounded corners=4pt, fill=UFSBlue, draw=UFSBlue] (2.15,0.5) rectangle (3.75,1.0);
        \node[font=\tiny, text=white] at (2.95,0.75) {Seguir};
      \end{tikzpicture}
      \end{center}
    \end{column}
  \end{columns}
\end{frame}

% ==================================================================
\section{Como um aplicativo chega ao celular}
% ==================================================================

\begin{frame}{Três caminhos para o mesmo aplicativo}
  \begin{center}
  \begin{tikzpicture}
    \node[cxg, minimum width=3.3cm] (n) at (-4.1,0)
      {\textbf{Nativo}\\ Kotlin/Java (Android)\\ Swift (iOS)};
    \node[cx,  minimum width=3.3cm] (m) at (0,0)
      {\textbf{Multiplataforma}\\ um código,\\ dois sistemas};
    \node[cxy, minimum width=3.3cm] (w) at (4.1,0)
      {\textbf{Web}\\ o site aberto\\ no navegador};
    \node[cel, minimum width=3.3cm, minimum height=0.9cm] (n2) at (-4.1,-1.6)
      {dois times, dois códigos,\\ dois cronogramas};
    \node[cel, minimum width=3.3cm, minimum height=0.9cm] (m2) at (0,-1.6)
      {um time, um código,\\ duas lojas};
    \node[cel, minimum width=3.3cm, minimum height=0.9cm] (w2) at (4.1,-1.6)
      {um endereço,\\ nenhuma loja};
    \draw[setag] (n) -- (n2);
    \draw[seta]  (m) -- (m2);
    \draw[setag] (w) -- (w2);
    \node[rot, align=center] at (0,-2.7)
      {A escolha não é sobre qual linguagem é mais bonita:\\
       é sobre custo de manutenção, acesso ao hardware e velocidade de entrega.};
  \end{tikzpicture}
  \end{center}
  \fonte{Eisenman (2017), cap.\ 1, ``Advantages of React Native''~\cite{eisenman2017learning}}
\end{frame}

\begin{frame}{Nativo: o teto mais alto e a conta mais cara}
  \begin{columns}[T]
    \begin{column}{0.5\textwidth}
      \textbf{O que você ganha}
      \begin{itemize}
        \item desempenho e fluidez máximos;
        \item acesso imediato a todo recurso novo do sistema;
        \item ferramentas maduras de depuração e perfil.
      \end{itemize}
    \end{column}
    \begin{column}{0.5\textwidth}
      \textbf{O que você paga}
      \begin{itemize}
        \item escrever \emph{duas vezes} a mesma tela;
        \item manter duas bases de código em sincronia;
        \item contratar duas especialidades distintas.
      \end{itemize}
    \end{column}
  \end{columns}
  \begin{block}{Quando compensa}
    Aplicativos que dependem de câmera, sensores ou gráficos em tempo real ---
    e equipes grandes o bastante para bancar dois times.
  \end{block}
  \fonte{Eisenman (2017), cap.\ 1~\cite{eisenman2017learning}}
\end{frame}

\begin{frame}{Web e híbrido: o mesmo HTML dentro de uma moldura}
  \begin{center}
  \begin{tikzpicture}
    \node[cxy, minimum width=3.6cm, minimum height=1.0cm] (app) at (0,0)
      {\textbf{Aplicativo híbrido}\\ instalado pela loja};
    \node[cxg, minimum width=3.0cm, minimum height=1.0cm] (wv) at (4.6,0)
      {\texttt{WebView}\\ (um navegador embutido)};
    \node[cel, minimum width=3.0cm, minimum height=0.8cm] (html) at (4.6,-1.5)
      {HTML + CSS + JS};
    \draw[seta] (app) -- (wv);
    \draw[setag] (wv) -- (html);
    \node[cel, minimum width=3.6cm, minimum height=1.3cm, align=left] at (0,-1.6)
      {A tela é uma página.\\ O sistema operacional\\ vê \emph{um} retângulo só.};
  \end{tikzpicture}
  \end{center}
  \begin{block}{A consequência que se sente no dedo}
    Rolagem, teclado e animações imitam o comportamento nativo em vez de
    \emph{serem} o comportamento nativo --- e a diferença aparece justamente
    nos gestos mais usados.
  \end{block}
  \fonte{Eisenman (2017), cap.\ 1, ``Hybrid Applications''~\cite{eisenman2017learning}}
\end{frame}

\begin{frame}{Pergunta}
  \begin{center}
    \vspace{1.2em}
    {\Large Quando você escreve \texttt{<Text>Olá</Text>} em React Native,}

    \vspace{0.6em}
    {\Large o que exatamente aparece na tela do celular?}

    \vspace{1.6em}
    \begin{tikzpicture}
      \node[cel, minimum width=3.2cm, minimum height=0.8cm] at (-3.4,0) {um \texttt{<p>} de HTML?};
      \node[cel, minimum width=3.2cm, minimum height=0.8cm] at (0,0) {um desenho em um \emph{canvas}?};
      \node[cel, minimum width=3.2cm, minimum height=0.8cm] at (3.4,0) {um componente do sistema?};
    \end{tikzpicture}
  \end{center}
\end{frame}

\begin{frame}{Resposta: um componente nativo de verdade}
  \begin{center}
  \begin{tikzpicture}
    \node[cx, minimum width=3.4cm, minimum height=0.9cm] (js) at (-3.6,0.9)
      {\textbf{Seu código}\\ JavaScript + JSX};
    \node[cxy, minimum width=3.4cm, minimum height=0.9cm] (rt) at (0.4,0.9)
      {\textbf{Tempo de execução}\\ do React Native};
    \node[cxr, minimum width=2.6cm, minimum height=0.9cm] (and) at (4.3,1.6)
      {\texttt{android.widget}\\ \texttt{.TextView}};
    \node[cxr, minimum width=2.6cm, minimum height=0.9cm] (ios) at (4.3,0.2)
      {\texttt{UIKit}\\ \texttt{.UITextView}};
    \draw[seta] (js) -- (rt);
    \draw[setar] (rt) -- (and);
    \draw[setar] (rt) -- (ios);
    \node[rot, align=center] at (0.4,-0.6) {a ponte traduz\\ a árvore de elementos};
    \node[cel, minimum width=9.6cm, minimum height=0.8cm, align=center] at (0.4,-1.7)
      {O \texttt{<Text>} \emph{não} é uma imitação: é o widget de texto do próprio sistema.\\
       Por isso a seleção de texto, a rolagem e a acessibilidade já vêm prontas.};
  \end{tikzpicture}
  \end{center}
  \fonte{React Native Docs, ``Core Components and Native Components''~\cite{reactnative2025docs}}
\end{frame}

\begin{frame}{Duas famílias de multiplataforma}
  \begin{center}
  \begin{tikzpicture}
    \node[cx, minimum width=4.4cm, minimum height=0.8cm] (a) at (-2.8,1.2)
      {\textbf{Usa os componentes do sistema}};
    \node[cxg, minimum width=4.4cm, minimum height=0.8cm] (b) at (2.8,1.2)
      {\textbf{Desenha tudo por conta própria}};
    \node[cel, minimum width=4.4cm, minimum height=1.5cm, align=left] (a2) at (-2.8,-0.2)
      {React Native\\[2pt]
       O botão é o botão do Android/iOS.\\
       Muda de aparência quando o\\ sistema muda.};
    \node[cel, minimum width=4.4cm, minimum height=1.5cm, align=left] (b2) at (2.8,-0.2)
      {Flutter\\[2pt]
       Pinta cada pixel em uma tela\\ própria. Aparência idêntica\\ nos dois sistemas.};
    \draw[seta] (a) -- (a2);
    \draw[setag] (b) -- (b2);
    \node[rot, align=center] at (0,-1.6)
      {Nenhuma das duas é ``a certa''. A pergunta é se você quer parecer\\
       com o sistema do usuário ou parecer sempre igual a si mesmo.};
  \end{tikzpicture}
  \end{center}
  \fonte{React Native Docs, ``Architecture Overview''~\cite{reactnative2025docs}}
\end{frame}

\begin{frame}{Recapitulando: o panorama}
  \begin{center}
  \small
  \begin{tabular}{@{}lccc@{}}
    \toprule
      & \textbf{Nativo} & \textbf{React Native} & \textbf{Web / híbrido} \\
    \midrule
    Bases de código        & duas       & uma            & uma \\
    O que desenha a tela   & o sistema  & o sistema      & o navegador \\
    Acesso a sensores      & total      & amplo          & limitado \\
    Custo de manutenção    & alto       & médio          & baixo \\
    Distribuição           & lojas      & lojas          & endereço web \\
    \bottomrule
  \end{tabular}
  \end{center}
  \begin{block}{A frase para levar}
    React Native escreve-se como uma página e executa-se como um aplicativo:
    a sua lógica é JavaScript, mas a interface é feita de componentes nativos.
  \end{block}
  \fonte{Eisenman (2017), cap.\ 1~\cite{eisenman2017learning}; React Native Docs~\cite{reactnative2025docs}}
\end{frame}

% ==================================================================
\section{JavaScript e React em trinta minutos}
% ==================================================================

\begin{frame}{Por que parar em JavaScript antes de abrir o editor}
  \begin{columns}[T]
    \begin{column}{0.5\textwidth}
      \textbf{O problema}\\[0.3em]
      Quase todo erro da primeira semana de React Native não é de React
      Native: é de JavaScript. Uma chave fora do lugar, uma função que
      devolve \texttt{undefined}, um \texttt{map} sem \texttt{return}.
    \end{column}
    \begin{column}{0.5\textwidth}
      \textbf{O recorte de hoje}\\[0.3em]
      Cinco construções, e só elas:
      \begin{itemize}
        \item \emph{arrow functions};
        \item desestruturação;
        \item \texttt{map};
        \item JSX;
        \item componentes e \emph{props}.
      \end{itemize}
    \end{column}
  \end{columns}
  \begin{block}{}
    A ideia central: em JavaScript, funções são valores ---
    dá para guardá-las em variáveis e passá-las como argumento.
  \end{block}
  \fonte{Flanagan (2020), cap.\ 8, ``Functions''~\cite{flanagan2020javascript}}
\end{frame}

\begin{frame}[fragile]{Arrow functions: a função como valor}
\begin{lstlisting}[style=js]
// forma tradicional: uma declaracao de funcao
function dobro(x) {
  return x * 2;
}

// a mesma funcao, escrita como valor guardado em uma constante
const dobro = (x) => x * 2;        // corpo curto: o return e' implicito

// quando o corpo tem mais de uma linha, volta o bloco e o return
const saudacao = (nome) => {
  const texto = 'Ola, ' + nome;
  return texto;
};
\end{lstlisting}
  \fonte{Flanagan (2020), \S 8.1.3, ``Arrow Functions''~\cite{flanagan2020javascript}}
\end{frame}

\begin{frame}[fragile]{Desestruturação: tirar campos de dentro de um objeto}
  \textbf{Desestruturação}: atribuição em que o lado esquerdo imita o formato
  do dado --- cada variável recebe o campo de mesmo nome (objeto) ou a mesma posição (array).
  \vspace{0.3em}
  \begin{columns}[T]
    \begin{column}{0.48\textwidth}
      \textbf{Sem desestruturação}
\begin{lstlisting}[style=jsx]
const aluno = { nome: 'Ada', nota: 9.5 };
const nome = aluno.nome;
const nota = aluno.nota;
\end{lstlisting}
    \end{column}
    \begin{column}{0.48\textwidth}
      \textbf{Com desestruturação}
\begin{lstlisting}[style=jsx]
const { nome, nota } = aluno;
// tambem funciona no parametro:
function resumo({ nome, nota }) {
  return nome + ': ' + nota;
}
// e em arrays, por posicao:
const [a, b] = [10, 20];
\end{lstlisting}
    \end{column}
  \end{columns}
  \begin{block}{}
    A forma de array é a que o \texttt{useState} usa: \texttt{const [valor, setValor] = useState(0)}.
  \end{block}
  \fonte{Flanagan (2020), \S 3.10.3, ``Destructuring Assignment'', e \S 8.3.5~\cite{flanagan2020javascript}}
\end{frame}

\begin{frame}[fragile]{\texttt{map}: de uma lista de dados para outra lista}
\begin{lstlisting}[style=js]
const habilidades = ['JavaScript', 'React Native', 'Flexbox'];

// map percorre a lista, aplica a funcao e devolve uma NOVA lista
const tamanhos = habilidades.map((h) => h.length);
// tamanhos vale [10, 12, 7]; habilidades continua intacta

const marcadas = habilidades.map((h) => '#' + h.toLowerCase());
// ['#javascript', '#react native', '#flexbox']
\end{lstlisting}
  \begin{block}{Por que isso importa em uma tela}
    Em React Native a lista de saída não é de textos: é de \emph{elementos
    de interface}. Uma etiqueta na tela para cada habilidade no vetor.
  \end{block}
  \fonte{Flanagan (2020), \S 7.8.1, ``\texttt{map()}''~\cite{flanagan2020javascript}}
\end{frame}

\begin{frame}[fragile]{Pergunta}
\begin{lstlisting}[style=js]
const notas = [7, 5, 9, 4];

const resultado = notas.map((n) => {
  n >= 6;
});
\end{lstlisting}
  \begin{center}
    \vspace{0.4em}
    {\large Quanto vale \texttt{resultado}?}\\[0.4em]
    {\small E o que seria preciso mudar para obter
     \texttt{[true, false, true, false]}?}
  \end{center}
\end{frame}

\begin{frame}[fragile]{Resposta: o \texttt{return} que faltou}
\begin{lstlisting}[style=js]
// resultado vale [undefined, undefined, undefined, undefined]
// o corpo em bloco { } executa a comparacao e joga fora o valor
\end{lstlisting}
  \begin{columns}[T]
    \begin{column}{0.48\textwidth}
\begin{lstlisting}[style=jsx]
// corrigindo com return
notas.map((n) => {
  return n >= 6;
});
\end{lstlisting}
    \end{column}
    \begin{column}{0.48\textwidth}
\begin{lstlisting}[style=jsx]
// ou com corpo curto
notas.map((n) => n >= 6);
\end{lstlisting}
    \end{column}
  \end{columns}
  \begin{alertblock}{A regra}
    Chave depois da seta significa bloco de comandos, e bloco de comandos
    exige \texttt{return}. Sem chave, o valor da expressão já é o retorno.
  \end{alertblock}
  \fonte{Flanagan (2020), \S 8.1.3~\cite{flanagan2020javascript}}
\end{frame}

\begin{frame}[fragile]{O que é JSX}
  \textbf{JSX} (\emph{JavaScript XML}) é uma \textbf{extensão de sintaxe} do
  JavaScript: tags no estilo XML/HTML, escritas no meio do código, que
  descrevem a tela. A forma do código imita a forma da interface.
  \begin{columns}[T]
    \begin{column}{0.47\textwidth}
      \small
      \begin{itemize}
        \item \textbf{Não é HTML}: as tags são componentes
          (\texttt{<View>}), não \texttt{<div>};
        \item \textbf{Não é texto}: é uma \emph{expressão}, um valor;
        \item \textbf{Não é JavaScript padrão}: o compilador
          Babel (embutido no Expo) traduz antes de rodar.
      \end{itemize}
    \end{column}
    \begin{column}{0.49\textwidth}
\begin{lstlisting}[style=js]
const nome = 'Ada';
// JSX e' um valor: vai numa variavel
const oi = <Text>Oi, {nome}</Text>;
// ... e num return
function Tela() {
  return <View>{oi}</View>;
}
\end{lstlisting}
    \end{column}
  \end{columns}
  \fonte{React Docs, ``Writing Markup with JSX''~\cite{react2025docs};
    Banks \& Porcello (2020), cap.\ 5~\cite{banks2020learning}}
\end{frame}

\begin{frame}{JSX por baixo: cada tag vira uma chamada de função}
  \begin{center}
  \begin{tikzpicture}[
      etapa/.style={rectangle, rounded corners=3pt, line width=0.8pt,
        minimum height=1.9cm, align=left, inner sep=5pt,
        font=\ttfamily\scriptsize, anchor=north west},
      num/.style={circle, text=white, font=\tiny\bfseries, inner sep=1.2pt,
        anchor=west},
      cab/.style={font=\scriptsize\bfseries, anchor=west},
    ]
    % 1. o que voce escreve
    \node[etapa, text width=2.6cm, draw=UFSBlue, fill=UFSBlue!8]
      (jsx) at (0,0) {<Text>\\ \ \ Olá\\ </Text>};
    % 2. o que o compilador gera
    \node[etapa, text width=3.5cm, draw=UFSYellow!80!black, fill=UFSYellow!20]
      (call) at ([xshift=0.9cm]jsx.north east)
      {React.createElement(\\ \ \ Text,\\ \ \ null,\\ \ \ 'Olá')};
    % 3. o que existe em execucao
    \node[etapa, text width=3.3cm, draw=UFSRed, fill=UFSRed!8]
      (el) at ([xshift=0.9cm]call.north east)
      {\{\\ \ \ type: Text,\\ \ \ props: \{\\ \ \ \ \ children: 'Olá' \}\\ \}};

    % cabecalhos numerados, todos na mesma linha
    \node[num, fill=UFSBlue] (n1) at ([yshift=0.3cm]jsx.north west) {1};
    \node[cab, text=UFSBlue] at (n1.east) {você escreve};
    \node[num, fill=UFSYellow!70!black] (n2) at ([yshift=0.3cm]call.north west) {2};
    \node[cab, text=UFSYellow!70!black] at (n2.east) {o compilador gera};
    \node[num, fill=UFSRed] (n3) at ([yshift=0.3cm]el.north west) {3};
    \node[cab, text=UFSRed] at (n3.east) {em execução vira};

    \draw[seta, line width=1pt] (jsx.east) -- (call.west);
    \draw[setar, line width=1pt] (call.east) -- (el.west);

    \node[rectangle, rounded corners=3pt, fill=black!5, text width=11.2cm,
          align=center, font=\small, inner sep=6pt, anchor=north]
      at ([yshift=-0.45cm]call.south)
      {JSX é \textbf{açúcar sintático} para chamadas de função.\\
       O celular nunca vê a marcação: vê o objeto que ela produz.};
  \end{tikzpicture}
  \end{center}
  \fonte{Banks \& Porcello (2020), cap.\ 5, ``React with JSX''~\cite{banks2020learning}}
\end{frame}

\begin{frame}[fragile]{Quatro regras de JSX que sempre pegam}
  \begin{columns}[T]
    \begin{column}{0.5\textwidth}
\begin{lstlisting}[style=jsx]
// 1. chaves entram em modo JavaScript
<Text>Total: {2 + 3}</Text>

// 2. so' pode haver UM elemento raiz
<View>
  <Text>a</Text>
  <Text>b</Text>
</View>

// 3. toda tag fecha, inclusive as vazias
<Image source={foto} />
\end{lstlisting}
    \end{column}
    \begin{column}{0.5\textwidth}
\begin{lstlisting}[style=jsx]
// 4. estilo e' objeto, nao string
<Text style={{ fontSize: 20 }}>
  Ola
</Text>

// as duas chaves: a de fora abre o
// modo JavaScript, a de dentro e'
// o literal de objeto
\end{lstlisting}
      \medskip
      \small
      Nomes de propriedade seguem JavaScript:\\
      \texttt{background-color} vira \texttt{backgroundColor}.
    \end{column}
  \end{columns}
  \fonte{Banks \& Porcello (2020), cap.\ 5~\cite{banks2020learning}}
\end{frame}

\begin{frame}[fragile]{Componente: uma função que devolve JSX}
\begin{lstlisting}[style=js]
// componente = funcao com nome MAIUSCULO que devolve JSX
function Saudacao() {
  return <Text>Bom dia!</Text>;
}

// usar o componente e' escrever uma tag com o nome dele
function App() {
  return (
    <View>
      <Saudacao />
      <Saudacao />
    </View>
  );
}
\end{lstlisting}
  \fonte{React Docs, ``Your First Component''~\cite{react2025docs}}
\end{frame}

\begin{frame}[fragile]{Props: os parâmetros do componente}
  \begin{columns}[T]
    \begin{column}{0.48\textwidth}
      \textbf{Quem recebe}
\begin{lstlisting}[style=jsx]
// as props chegam como UM objeto;
// a desestruturacao abre esse objeto
function Saudacao({ nome, horario }) {
  return <Text>Bom {horario}, {nome}!</Text>;
}
\end{lstlisting}
    \end{column}
    \begin{column}{0.48\textwidth}
      \textbf{Quem envia}
\begin{lstlisting}[style=jsx]
<View>
  <Saudacao nome="Ada" horario="dia" />
  <Saudacao nome="Alan" horario="tarde" />
</View>
\end{lstlisting}
    \end{column}
  \end{columns}
  \begin{alertblock}{Props são somente leitura}
    Quem decide o valor é quem chamou o componente. Para mudar algo de
    dentro, é preciso \emph{estado} --- o último bloco da aula.
  \end{alertblock}
  \fonte{React Docs, ``Passing Props to a Component''~\cite{react2025docs}}
\end{frame}

\begin{frame}{Recapitulando: JavaScript e React}
  \begin{center}
  \begin{tikzpicture}
    \node[cel, minimum width=11.0cm, minimum height=0.55cm] (a) at (0,0)
      {\texttt{(x) => x * 2} \quad --- função como valor; sem chave, o retorno é implícito};
    \node[cel, minimum width=11.0cm, minimum height=0.55cm] (b) at (0,-0.75)
      {\texttt{const \{ nome \} = aluno} \quad --- abre o objeto; é assim que props chegam};
    \node[cel, minimum width=11.0cm, minimum height=0.55cm] (c) at (0,-1.5)
      {\texttt{lista.map(...)} \quad --- uma lista de dados vira uma lista de elementos};
    \node[cel, minimum width=11.0cm, minimum height=0.55cm] (d) at (0,-2.25)
      {\texttt{<View><Text/></View>} \quad --- JSX: um elemento raiz, tudo fechado};
    \node[cel, minimum width=11.0cm, minimum height=0.55cm] (e) at (0,-3.0)
      {\texttt{function Cartao(\{ nome \})} \quad --- componente com nome maiúsculo e props};
  \end{tikzpicture}
  \end{center}
\end{frame}

% ==================================================================
\section{A primeira tela no seu celular}
% ==================================================================

\begin{frame}{Por que Snack e Expo Go, e não a instalação completa}
  \begin{center}
  \begin{tikzpicture}
    \node[cx, minimum width=3.2cm, minimum height=0.9cm] (nav) at (-4.0,0)
      {\textbf{Navegador}\\ \texttt{snack.expo.dev}};
    \node[cxy, minimum width=3.2cm, minimum height=0.9cm] (srv) at (0,0)
      {\textbf{Servidor do Expo}\\ empacota o seu código};
    \node[cxr, minimum width=3.2cm, minimum height=0.9cm] (cel) at (4.0,0)
      {\textbf{Expo Go}\\ no seu celular};
    \draw[seta] (nav) -- node[rot, above] {salva} (srv);
    \draw[setar] (srv) -- node[rot, above] {QR code} (cel);
    \draw[setag] (cel.south) to[bend left=25] node[rot, below] {você edita e a tela recarrega sozinha} (nav.south);
  \end{tikzpicture}
  \end{center}
  \begin{block}{O que isso nos poupa hoje}
    Android Studio, Xcode, SDK, emulador e JDK --- somados, várias horas de
    instalação e um disco cheio. Nada disso é necessário para as ideias de hoje.
  \end{block}
  \fonte{Expo Docs, ``Snack'' e ``Expo Go''~\cite{expo2025docs}}
\end{frame}

\begin{frame}{Cinco passos, cinco minutos}
  \begin{enumerate}
    \item No celular, instale o \textbf{Expo Go} (Play Store ou App Store).
    \item No computador, abra \texttt{snack.expo.dev} em uma aba nova.
    \item Na barra direita, escolha \textbf{My Device}: aparece um QR code.
    \item Abra o Expo Go e leia o QR code (no iPhone, use a câmera).
    \item Apague o conteúdo de \texttt{App.js} e cole o código do próximo slide.
  \end{enumerate}
  \begin{alertblock}{Se a rede da universidade bloquear a conexão}
    Use o botão \textbf{Web} do próprio Snack: a tela roda no navegador.
    Perde-se o celular de verdade, mas nenhum conceito da aula.
  \end{alertblock}
  \fonte{Expo Docs, ``Get started''~\cite{expo2025docs}}
\end{frame}

\begin{frame}[fragile]{O ``Olá mundo'' que vamos rodar}
\begin{lstlisting}[style=js]
import { View, Text, StyleSheet } from 'react-native';

export default function App() {
  return (
    <View style={estilos.tela}>
      <Text style={estilos.titulo}>Ola, COMP0511!</Text>
    </View>
  );
}

const estilos = StyleSheet.create({
  tela: { flex: 1, justifyContent: 'center', alignItems: 'center' },
  titulo: { fontSize: 28, fontWeight: 'bold' },
});
\end{lstlisting}
  \fonte{React Native Docs, ``Core Components''~\cite{reactnative2025docs}}
\end{frame}

\begin{frame}{Faça agora --- 5 minutos}
  \begin{block}{Tarefa 1}
    \begin{enumerate}
      \item Ponha o ``Olá mundo'' para rodar no seu aparelho.
      \item Troque o texto pelo seu próprio nome.
      \item Mude \texttt{fontSize} para \texttt{48} e veja a tela recarregar sozinha.
      \item Apague a linha \texttt{flex: 1} e observe para onde o texto vai.
    \end{enumerate}
  \end{block}
  \begin{columns}[T]
    \begin{column}{0.5\textwidth}
      \textbf{Se travar}\\
      Chame; quem já conseguiu ajuda o vizinho. Ninguém segue para o próximo
      bloco sem ter visto o próprio nome na tela.
    \end{column}
    \begin{column}{0.5\textwidth}
      \textbf{Guarde o endereço}\\
      O Snack gera uma URL. Salve-a: é nela que o cartão de perfil vai crescer
      pelo resto da aula.
    \end{column}
  \end{columns}
\end{frame}

\begin{frame}{Quando não funciona: os quatro suspeitos de sempre}
  \begin{center}
  \small
  \begin{tabular}{@{}p{4.6cm}p{7.0cm}@{}}
    \toprule
    \textbf{Sintoma} & \textbf{Causa provável e o que fazer} \\
    \midrule
    O QR code não conecta
      & Celular e computador em redes diferentes. Use \texttt{Web} ou troque a rede. \\
    Tela branca, sem erro
      & Faltou o \texttt{export default} antes de \texttt{function App}. \\
    \texttt{Text strings must be rendered within a <Text>}
      & Texto solto dentro de \texttt{<View>}. Envolva-o em \texttt{<Text>}. \\
    A imagem não aparece
      & \texttt{Image} com \texttt{uri} precisa de \texttt{width} e \texttt{height} no estilo. \\
    \bottomrule
  \end{tabular}
  \end{center}
  \fonte{React Native Docs, ``Troubleshooting''~\cite{reactnative2025docs}}
\end{frame}

% ==================================================================
\section{Componentes, estilos e layout}
% ==================================================================

\begin{frame}{O catálogo mínimo: cinco componentes bastam hoje}
  \begin{center}
  \small
  \begin{tabular}{@{}llp{6.6cm}@{}}
    \toprule
    \textbf{Componente} & \textbf{Equivale a} & \textbf{Para que serve} \\
    \midrule
    \texttt{View}      & \texttt{<div>}    & caixa: agrupa, posiciona, recebe cor e borda \\
    \texttt{Text}      & \texttt{<p>}      & todo texto visível; nada de texto solto \\
    \texttt{Image}     & \texttt{<img>}    & imagem local ou vinda de uma URL \\
    \texttt{TextInput} & \texttt{<input>}  & entrada de texto do usuário \\
    \texttt{Pressable} & \texttt{<button>} & qualquer área que reage ao toque \\
    \bottomrule
  \end{tabular}
  \end{center}
  \begin{block}{A regra que o iniciante mais esquece}
    Texto só existe dentro de \texttt{<Text>}. Escrever palavras diretamente
    dentro de uma \texttt{<View>} é erro de execução, não aviso.
  \end{block}
  \fonte{React Native Docs, ``Core Components and APIs''~\cite{reactnative2025docs}}
\end{frame}

\begin{frame}[fragile]{Os cinco em uma tela só}
\begin{lstlisting}[style=js]
<View style={estilos.tela}>
  <Text style={estilos.titulo}>Catalogo</Text>

  <Image style={estilos.foto}
         source={{ uri: 'https://picsum.photos/id/1005/200/200' }} />

  <TextInput style={estilos.campo} placeholder="Digite seu nome" />

  <Pressable style={estilos.botao} onPress={() => console.log('toquei')}>
    <Text style={estilos.textoBotao}>Tocar aqui</Text>
  </Pressable>
</View>
\end{lstlisting}
  \begin{block}{}
    \texttt{Pressable} não desenha nada: ele envolve o que deve ser tocável.
  \end{block}
  \fonte{React Native Docs, ``Pressable'' e ``Image''~\cite{reactnative2025docs}}
\end{frame}

\begin{frame}[fragile]{\texttt{StyleSheet.create}: estilo separado da estrutura}
  \begin{columns}[T]
    \begin{column}{0.48\textwidth}
      \textbf{Inline: rápido e bagunçado}
\begin{lstlisting}[style=jsx]
<Text style={{ fontSize: 20,
               color: 'blue' }}>
  Ola
</Text>
\end{lstlisting}
      \small Um objeto novo a cada renderização; a estrutura some no meio da
      formatação.
    \end{column}
    \begin{column}{0.48\textwidth}
      \textbf{Com \texttt{StyleSheet}}
\begin{lstlisting}[style=jsx]
<Text style={estilos.titulo}>Ola</Text>

const estilos = StyleSheet.create({
  titulo: { fontSize: 20,
            color: 'blue' },
});
\end{lstlisting}
      \small Nomes com significado, reaproveitáveis, e o JSX volta a se ler
      como estrutura.
    \end{column}
  \end{columns}
  \begin{block}{Combinando estilos}
    \texttt{style=\{[estilos.botao, estilos.botaoAtivo]\}} --- o vetor aplica
    da esquerda para a direita; o último vence nos campos repetidos.
  \end{block}
  \fonte{React Native Docs, ``Style'' e ``StyleSheet''~\cite{reactnative2025docs}}
\end{frame}

\begin{frame}[fragile]{Pergunta}
\begin{lstlisting}[style=js]
const estilos = StyleSheet.create({
  tela: { justifyContent: 'center', alignItems: 'center' },
});
\end{lstlisting}
  \begin{center}
    \vspace{0.5em}
    {\large A tela tem um único \texttt{<Text>} e o estilo acima.}\\[0.5em]
    {\large Por que o texto continua colado no topo?}
  \end{center}
\end{frame}

\begin{frame}[fragile]{Resposta: centralizar dentro de quê?}
\begin{lstlisting}[style=js]
tela: { flex: 1, justifyContent: 'center', alignItems: 'center' }
\end{lstlisting}
  \begin{center}
  \begin{tikzpicture}[scale=0.9, transform shape]
    \draw[draw=UFSGray, fill=black!3] (0,0) rectangle (3.4,2.6);
    \draw[draw=UFSRed, dashed] (0.1,2.1) rectangle (3.3,2.5);
    \node[font=\tiny, text=UFSRed] at (1.7,2.3) {a View tem só a altura do texto};
    \node[font=\tiny, text=UFSGray] at (1.7,0.5) {sem \texttt{flex: 1}};
    \draw[draw=UFSGray, fill=black!3] (5.0,0) rectangle (8.4,2.6);
    \draw[draw=UFSBlue, dashed] (5.1,0.1) rectangle (8.3,2.5);
    \node[font=\tiny, text=UFSBlue] at (6.7,1.3) {a View ocupa a tela inteira};
    \node[font=\tiny, text=UFSGray] at (6.7,-0.35) {com \texttt{flex: 1}};
  \end{tikzpicture}
  \end{center}
  \begin{alertblock}{}
    \texttt{justifyContent} distribui o conteúdo \emph{no espaço que a caixa tem}.
    Se a caixa encolheu até caber no conteúdo, não há espaço para distribuir.
  \end{alertblock}
  \fonte{React Native Docs, ``Layout with Flexbox''~\cite{reactnative2025docs}}
\end{frame}

\begin{frame}{Flexbox: três propriedades resolvem quase tudo}
  \begin{center}
  \begin{tikzpicture}[scale=0.92, transform shape]
    % coluna
    \node[font=\scriptsize\bfseries] at (1.5,3.1) {flexDirection: 'column'};
    \node[font=\tiny, text=UFSGray] at (1.5,2.8) {(o padrão no React Native)};
    \draw[draw=UFSGray] (0,0) rectangle (3.0,2.6);
    \draw[fill=UFSBlue!25, draw=UFSBlue] (0.4,1.85) rectangle (2.6,2.35);
    \draw[fill=UFSBlue!25, draw=UFSBlue] (0.4,1.15) rectangle (2.6,1.65);
    \draw[fill=UFSBlue!25, draw=UFSBlue] (0.4,0.45) rectangle (2.6,0.95);
    \draw[setar] (-0.35,2.5) -- (-0.35,0.1);
    \node[rot, rotate=90, text=UFSRed] at (-0.6,1.3) {eixo principal};
    \draw[setag] (0.1,-0.3) -- (2.9,-0.3);
    \node[rot] at (1.5,-0.55) {eixo cruzado};

    % linha
    \node[font=\scriptsize\bfseries] at (7.5,3.1) {flexDirection: 'row'};
    \draw[draw=UFSGray] (6.0,0) rectangle (9.0,2.6);
    \draw[fill=UFSYellow!35, draw=UFSYellow!80!black] (6.3,0.7) rectangle (7.0,1.9);
    \draw[fill=UFSYellow!35, draw=UFSYellow!80!black] (7.15,0.7) rectangle (7.85,1.9);
    \draw[fill=UFSYellow!35, draw=UFSYellow!80!black] (8.0,0.7) rectangle (8.7,1.9);
    \draw[setar] (6.1,-0.3) -- (8.9,-0.3);
    \node[rot, text=UFSRed] at (7.5,-0.55) {eixo principal};

    % legenda
    \node[cel, minimum width=4.0cm, minimum height=1.6cm, align=left] at (12.0,1.3)
      {\texttt{flexDirection}\\ escolhe o eixo principal\\[4pt]
       \texttt{justifyContent}\\ distribui no eixo principal\\[4pt]
       \texttt{alignItems}\\ alinha no eixo cruzado};
  \end{tikzpicture}
  \end{center}
  \fonte{W3C, CSS Flexible Box Layout L1, \S 5--6~\cite{w3c2018flexbox}; React Native Docs~\cite{reactnative2025docs}}
\end{frame}

\begin{frame}{Flexbox no celular não é igual ao da web}
  \begin{center}
  \small
  \begin{tabular}{@{}lll@{}}
    \toprule
      & \textbf{CSS na web} & \textbf{React Native} \\
    \midrule
    \texttt{flexDirection} padrão & \texttt{row}                & \texttt{column} \\
    \texttt{flex}                 & três valores               & um número só \\
    Unidades                      & \texttt{px}, \texttt{em}, \texttt{\%} & número puro (pontos) \\
    Herança de fonte              & filhos herdam              & cada \texttt{Text} tem o seu \\
    Nome das propriedades         & \texttt{background-color}   & \texttt{backgroundColor} \\
    \bottomrule
  \end{tabular}
  \end{center}
  \begin{block}{Por que \texttt{column} é o padrão}
    Telas de celular são altas e estreitas: empilhar é o caso comum.
    A web herdou do texto, que corre na horizontal.
  \end{block}
  \fonte{React Native Docs, ``Layout with Flexbox''~\cite{reactnative2025docs};
         Eisenman (2017), cap.\ 5~\cite{eisenman2017learning}}
\end{frame}

\begin{frame}[fragile]{Exemplo integrado: o cabeçalho do cartão}
  \begin{columns}[T]
    \begin{column}{0.56\textwidth}
\begin{lstlisting}[style=jsx]
<View style={estilos.cabecalho}>
  <Image style={estilos.foto}
         source={{ uri: PERFIL.foto }} />
  <View style={estilos.identidade}>
    <Text style={estilos.nome}>Ada Lovelace</Text>
    <Text style={estilos.usuario}>@ada</Text>
  </View>
</View>

cabecalho: { flexDirection: 'row',
             alignItems: 'center' },
foto:      { width: 72, height: 72,
             borderRadius: 36 },
identidade:{ marginLeft: 14 },
\end{lstlisting}
    \end{column}
    \begin{column}{0.40\textwidth}
      \begin{center}
      \begin{tikzpicture}
        \draw[draw=UFSGray, dashed] (0,0) rectangle (4.2,1.5);
        \node[rot, anchor=south west] at (0,1.55) {cabecalho: row};
        \draw[fill=UFSBlue!25, draw=UFSBlue] (0.75,0.75) circle (0.5);
        \draw[draw=UFSGray, dashed] (1.5,0.35) rectangle (4.05,1.15);
        \node[rot, anchor=west] at (1.6,0.95) {Ada Lovelace};
        \node[rot, anchor=west] at (1.6,0.6) {@ada};
        \node[rot, anchor=north west] at (1.5,0.3) {identidade: column};
      \end{tikzpicture}
      \end{center}
      \medskip
      \small Uma \texttt{View} em linha, contendo uma \texttt{View} em coluna:
      é assim que quase todo cartão de rede social é montado.
    \end{column}
  \end{columns}
\end{frame}

\begin{frame}{Recapitulando: componentes e layout}
  \begin{block}{O que aprendemos}
    \begin{itemize}
      \item cinco componentes cobrem a tela de hoje, e texto só vive dentro de \texttt{Text};
      \item \texttt{StyleSheet.create} devolve o JSX à sua função de descrever estrutura;
      \item \texttt{flex: 1} é o que dá à caixa o espaço que ela vai distribuir;
      \item \texttt{flexDirection} escolhe o eixo; \texttt{justifyContent} e
        \texttt{alignItems} posicionam nele e no eixo cruzado.
    \end{itemize}
  \end{block}
  \begin{alertblock}{Ainda falta uma coisa}
    Tudo o que fizemos até aqui é uma fotografia: bonita e imóvel.
    O botão não conta, o campo não lembra.
  \end{alertblock}
\end{frame}

% ==================================================================
\section{Telas que reagem: \texttt{useState}}
% ==================================================================

\begin{frame}[fragile]{Pergunta}
\begin{lstlisting}[style=js]
export default function App() {
  let contagem = 0;

  return (
    <View>
      <Text>{contagem}</Text>
      <Pressable onPress={() => { contagem = contagem + 1; }}>
        <Text>somar 1</Text>
      </Pressable>
    </View>
  );
}
\end{lstlisting}
  \begin{center}
    {\large Você toca cinco vezes no botão. O que a tela mostra?}
  \end{center}
\end{frame}

\begin{frame}{Resposta: zero, sempre zero}
  \begin{center}
  \begin{tikzpicture}
    \node[cx, minimum width=2.9cm, minimum height=0.85cm] (r) at (-3.6,0.9) {a função \texttt{App}\\ é executada};
    \node[cxy, minimum width=2.9cm, minimum height=0.85cm] (t) at (0,0.9) {a tela é desenhada\\ com \texttt{contagem = 0}};
    \node[cxr, minimum width=2.9cm, minimum height=0.85cm] (e) at (3.6,0.9) {você toca:\\ \texttt{contagem} vira 1};
    \draw[seta] (r) -- (t);
    \draw[seta] (t) -- (e);
    \node[cel, minimum width=10.0cm, minimum height=1.0cm, align=center] at (0,-0.6)
      {A variável mudou na memória --- e ninguém pediu para redesenhar a tela.\\
       Pior: se algo redesenhar, a função roda de novo e \texttt{contagem} volta a valer 0.};
    \node[rot, align=center] at (0,-1.6)
      {Falta um mecanismo que faça duas coisas ao mesmo tempo:\\
       \emph{lembrar} o valor entre as renderizações e \emph{avisar} que a tela mudou.};
  \end{tikzpicture}
  \end{center}
  \fonte{React Docs, ``State: A Component's Memory''~\cite{react2025docs}}
\end{frame}

\begin{frame}[fragile]{\texttt{useState}: memória e aviso na mesma linha}
\begin{lstlisting}[style=js]
import { useState } from 'react';

const [contagem, setContagem] = useState(0);
//     |            |                    +-- valor inicial (so' na 1a vez)
//     |            +----------- troca o valor E pede um novo desenho
//     +------------------------ valor atual, valido nesta renderizacao
\end{lstlisting}
  \begin{center}
  \begin{tikzpicture}
    \node[cx, minimum width=2.6cm, minimum height=0.7cm] (a) at (-3.9,0) {estado};
    \node[cxy, minimum width=2.6cm, minimum height=0.7cm] (b) at (-0.6,0) {tela desenhada};
    \node[cxr, minimum width=2.6cm, minimum height=0.7cm] (c) at (2.7,0) {toque do usuário};
    \node[cx, minimum width=2.6cm, minimum height=0.7cm] (d) at (6.0,0) {\texttt{setContagem}};
    \draw[seta] (a) -- (b);
    \draw[seta] (b) -- (c);
    \draw[seta] (c) -- (d);
    \draw[seta] (d.north) to[bend right=30] (a.north);
  \end{tikzpicture}
  \end{center}
  \fonte{React Docs, ``State: A Component's Memory''~\cite{react2025docs};
         Banks \& Porcello (2020), cap.\ 6~\cite{banks2020learning}}
\end{frame}

\begin{frame}[fragile]{Contador: o programa inteiro}
\begin{lstlisting}[style=js]
import { useState } from 'react';
import { View, Text, Pressable, StyleSheet } from 'react-native';

export default function App() {
  const [contagem, setContagem] = useState(0);

  return (
    <View style={estilos.tela}>
      <Text style={estilos.numero}>{contagem}</Text>
      <Pressable style={estilos.botao} onPress={() => setContagem(contagem + 1)}>
        <Text style={estilos.textoBotao}>+1</Text>
      </Pressable>
    </View>
  );
}
\end{lstlisting}
  \fonte{React Native Docs, ``Using Hooks''~\cite{reactnative2025docs}}
\end{frame}

\begin{frame}[fragile]{Campo controlado: o estado manda no que aparece}
  \begin{columns}[T]
    \begin{column}{0.54\textwidth}
\begin{lstlisting}[style=jsx]
const [nome, setNome] = useState('');

<TextInput
  value={nome}
  onChangeText={setNome}
  placeholder="digite aqui"
/>

<Text>
  {nome.length === 0
    ? 'Ainda nao digitou nada.'
    : 'Ola, ' + nome + '!'}
</Text>
\end{lstlisting}
    \end{column}
    \begin{column}{0.42\textwidth}
      \begin{center}
      \begin{tikzpicture}
        \node[cel, minimum width=2.4cm, minimum height=0.6cm] (s) at (0,1.5) {estado \texttt{nome}};
        \node[cel, minimum width=2.4cm, minimum height=0.6cm] (c) at (0,0) {o que o campo exibe};
        \node[cel, minimum width=2.4cm, minimum height=0.6cm] (k) at (0,-1.5) {\texttt{onChangeText}};
        \draw[seta] (s) -- node[rot, right] {\texttt{value}} (c);
        \draw[setar] (c) -- node[rot, right] {você digita} (k);
        \draw[seta] (k.west) -- ++(-0.45,0) |- node[rot, left, pos=0.25] {\texttt{setNome}} (s.west);
      \end{tikzpicture}
      \end{center}
      \small O campo não guarda nada: mostra o estado e avisa a mudança.
    \end{column}
  \end{columns}
  \fonte{React Native Docs, ``Handling Text Input''~\cite{reactnative2025docs}}
\end{frame}

\begin{frame}[fragile]{Três armadilhas do estado}
  \begin{block}{1. \texttt{set} não é atribuição imediata}
\begin{lstlisting}[style=jsx]
setContagem(contagem + 1);
console.log(contagem);  // ainda o valor ANTIGO
\end{lstlisting}
  \end{block}
  \begin{block}{2. Se o novo valor depende do antigo, use a forma funcional}
\begin{lstlisting}[style=jsx]
setContagem((n) => n + 1);
\end{lstlisting}
  \end{block}
  \begin{block}{3. Nunca altere o valor por dentro}
\begin{lstlisting}[style=jsx]
lista.push('novo');            // ERRADO: React nao percebe
setLista([...lista, 'novo']);  // CERTO: um vetor novo
\end{lstlisting}
  \end{block}
  \fonte{React Docs, ``Updating Arrays in State''~\cite{react2025docs}}
\end{frame}

\begin{frame}{Tarefa 2 --- o cartão de perfil é seu}
  \begin{columns}[T]
    \begin{column}{0.56\textwidth}
      \textbf{O que construir}
      \begin{itemize}
        \item cabeçalho com foto redonda, nome e usuário;
        \item bio editável por um campo controlado;
        \item etiquetas de habilidades geradas com \texttt{map};
        \item botão de curtir que conta;
        \item botão Seguir/Seguindo que troca texto e cor.
      \end{itemize}
      \textbf{Como trabalhar:} em duplas, no mesmo Snack, a partir do
      esqueleto.
    \end{column}
    \begin{column}{0.42\textwidth}
      \begin{center}
      \begin{tikzpicture}
        \node[cx, minimum width=4.4cm, minimum height=0.62cm] (c1) at (0,0)
          {\textbf{1.} layout, Flexbox e \texttt{map}};
        \node[cx, minimum width=4.4cm, minimum height=0.62cm] (c2) at (0,-0.85)
          {\textbf{2.} contador de curtidas};
        \node[cx, minimum width=4.4cm, minimum height=0.62cm] (c3) at (0,-1.70)
          {\textbf{3.} bio controlada};
        \node[cx, minimum width=4.4cm, minimum height=0.62cm] (c4) at (0,-2.55)
          {\textbf{4.} Seguir / Seguindo};
        \draw[seta] (c1) -- (c2);
        \draw[seta] (c2) -- (c3);
        \draw[seta] (c3) -- (c4);
      \end{tikzpicture}
      \end{center}
    \end{column}
  \end{columns}
\end{frame}

\begin{frame}{O que cada checkpoint cobra}
  \begin{center}
  \small
  \begin{tabular}{@{}clp{6.4cm}@{}}
    \toprule
    \textbf{\#} & \textbf{Entrega} & \textbf{Conceito que ele verifica} \\
    \midrule
    1 & cartão estático completo
      & \texttt{flexDirection: 'row'}; \texttt{map} com \texttt{key} \\
    2 & botão de curtir contando
      & \texttt{useState} numérico e a forma funcional do \texttt{set} \\
    3 & bio que muda enquanto se digita
      & \texttt{value} + \texttt{onChangeText} no mesmo estado \\
    4 & Seguir vira Seguindo
      & estado booleano decidindo texto, estilo e um valor derivado \\
    \bottomrule
  \end{tabular}
  \end{center}
  \begin{alertblock}{Checkpoint 4, a pegadinha boa}
    O total de seguidores \emph{não} é estado novo: calcula-se a partir de
    \texttt{seguindo}. Estado a mais é fonte de inconsistência.
  \end{alertblock}
  \fonte{React Docs, ``Choosing the State Structure''~\cite{react2025docs}}
\end{frame}

% ==================================================================
\section{Fechamento}
% ==================================================================

\begin{frame}{O caminho que percorremos}
  \begin{center}
  \begin{tikzpicture}
    \node[cel, minimum width=11.0cm, minimum height=0.55cm] at (0,0)
      {React Native escreve-se como uma página e executa componentes \textbf{nativos}};
    \node[cel, minimum width=11.0cm, minimum height=0.55cm] at (0,-0.75)
      {arrow functions, desestruturação e \texttt{map} --- o JavaScript que a tela exige};
    \node[cel, minimum width=11.0cm, minimum height=0.55cm] at (0,-1.5)
      {JSX vira chamada de função; componente é função com nome maiúsculo};
    \node[cel, minimum width=11.0cm, minimum height=0.55cm] at (0,-2.25)
      {\texttt{flex: 1} dá o espaço; \texttt{flexDirection} escolhe o eixo};
    \node[cel, minimum width=11.0cm, minimum height=0.55cm] at (0,-3.0)
      {\texttt{useState} é memória \emph{e} aviso: sem ele a tela não se redesenha};
  \end{tikzpicture}
  \end{center}
\end{frame}

\begin{frame}{Para casa}
  \begin{block}{Tarefa}
    Transforme o cartão de perfil em um \textbf{cartão de contato} com uma
    lista de contatos:
    \begin{enumerate}
      \item extraia um componente \texttt{Cartao} que receba
        \texttt{nome}, \texttt{usuario}, \texttt{foto} e \texttt{habilidades} por props;
      \item monte um vetor com \textbf{quatro} pessoas e renderize os quatro
        cartões com \texttt{map};
      \item cada cartão mantém o próprio contador de curtidas
        (curtir um não pode afetar os outros);
      \item acrescente um \texttt{TextInput} no topo que filtra os cartões pelo nome.
    \end{enumerate}
  \end{block}
  \textbf{Dica do item 3:} onde deve morar o estado das curtidas ---
  no \texttt{App} ou dentro do \texttt{Cartao}?
\end{frame}

\begin{frame}{Na próxima aula}
  \begin{columns}[T]
    \begin{column}{0.5\textwidth}
      \textbf{De onde partimos}\\[0.3em]
      Uma tela só, com os dados escritos no próprio arquivo. Quatro cartões
      já deixam a tela sem espaço para rolar.
    \end{column}
    \begin{column}{0.5\textwidth}
      \textbf{Para onde vamos}\\[0.3em]
      \begin{itemize}
        \item listas longas com \texttt{FlatList};
        \item \texttt{ScrollView} e áreas seguras;
        \item navegação entre telas;
        \item dados vindos da rede.
      \end{itemize}
    \end{column}
  \end{columns}
  \begin{block}{Antes da próxima aula}
    Deixe o Expo Go instalado e o Snack de hoje salvo --- vamos continuar
    exatamente deste ponto.
  \end{block}
\end{frame}

\section*{Referências}
\begin{frame}[allowframebreaks]{Referências}
  \nocite{*}
  \bibliographystyle{plain}
  \bibliography{referencias}
\end{frame}

\end{document}
